% Generated draft: requires mathematical and source review.
\documentclass[11pt]{article}
\usepackage[margin=25mm]{geometry}
\usepackage{fontspec}
\setmainfont{Latin Modern Roman}
\usepackage{amsmath,amssymb,mathtools,unicode-math}
\setmathfont{Latin Modern Math}
\usepackage{enumitem,needspace,xurl,hyperref}
\setlength{\parindent}{0pt}
\setlength{\parskip}{5pt}
\setlength{\emergencystretch}{3em}
\newcommand{\N}{\mathbb{N}}
\newcommand{\Z}{\mathbb{Z}}
\newcommand{\Q}{\mathbb{Q}}
\newcommand{\R}{\mathbb{R}}
\newcommand{\C}{\mathbb{C}}
\newcommand{\F}{\mathbb{F}}
\newcommand{\A}{\mathbb{A}}
\newcommand{\PP}{\mathbb{P}}
\newcommand{\E}{\mathbb{E}}
\newcommand{\eps}{\varepsilon}
\newcommand{\dd}{\,\mathrm{d}}
\begin{document}
\section*{UnsolvedMath ID: 800484}
\textbf{Source problem code:} catalogue entry
\subsection{Definitions and mathematical statement}
Let \(\Omega\subset\R^n\) be a domain with smooth boundary, and let \(P\in\Omega\). The source does not specify the dimension, boundedness of \(\Omega\), or the meaning of a generic choice of \(P\). Normalize the common speed of the photons to one. Between reflections a photon follows a straight line; at a regular impact point \(x\in\partial\Omega\), its incoming and outgoing unit velocities satisfy the specular reflection law
\[
 v_{\mathrm{out}}=v_{\mathrm{in}}-2\langle v_{\mathrm{in}},\nu(x)\rangle\nu(x),
\]
where \(\nu(x)\) is the outward unit normal.

For \(t\geq0\), let \(W_t(P)\) denote the set of positions at time \(t\) of photons starting at \(P\) in all unit directions and following this reflection law. This definition requires a convention for trajectories with tangential impacts or other failures of continuation up to time \(t\); the source supplies no such convention. Restricting to trajectories defined through time \(t\) is one possible interpretation, but is not specified by the source.

\textbf{Problem (Serge Tabachnikov, Problem 6).} For a generic interior point \(P\), does the reflected wave front \(W_t(P)\) become dense as \(t\to\infty\)?

The phrase ``become dense'' is not defined in the supplied record. A precise possible interpretation is eventual uniform spatial approximation:
\[
 \forall\varepsilon>0\ \exists T\geq0\ \forall t\geq T\ \forall x\in\Omega\
 \exists y\in W_t(P):\quad \|x-y\|<\varepsilon.
\tag{1}
\]
A different, weaker interpretation is eventual intersection with each fixed nonempty relatively open subset:
\[
 \forall U\subseteq\Omega\ \bigl(U\text{ nonempty and relatively open}\bigr)
 \quad\exists T_U\geq0\ \forall t\geq T_U:\quad W_t(P)\cap U\neq\varnothing.
\tag{2}
\]
The supplied wording does not determine which interpretation is intended, or whether approximation only along a sequence of times is intended. Formula (1) also requires attention to boundedness: at any finite time every unit-speed photon lies within distance \(t\) of \(P\), so an unbounded domain cannot satisfy (1). No boundedness hypothesis is stated in the record.

The source also mentions an analogous question for smooth Riemannian manifolds, without specifying the class of manifolds, boundary conditions, or the precise limiting-density property. It reports that even the case described as a ``round circle'' is open; this is a supplied status qualification, not an independently verified mathematical assessment.
\subsection{Short English statement}
Serge Tabachnikov asks whether light emitted in every direction from a generic interior point eventually produces reflected wave fronts that approximate the entire billiard domain. The meaning of genericity, the limiting-density condition, the domain hypotheses, and exceptional-trajectory conventions remain unspecified. The source also mentions a smooth Riemannian analogue and reports the round-circle case as open.
\subsection{Sources}
\begin{itemize}
\item Local UnsolvedMath registry, supplied complete record: permanent ID \(800484\), exact source problem code \texttt{catalogue entry}; title: \emph{Billiards and Quantitative Symplectic Geometry --- Serge Tabachnikov --- Problem 6}. The registry labels the problem \texttt{open}; that label is supplied source data and has not been independently verified.
\item Albach, Barraud et al., \emph{Open Problems in Billiards and Quantitative Symplectic Geometry} (2026), Serge Tabachnikov, Problem 6, \url{https://arxiv.org/abs/2602.12896}. The supplied background records access on 2026-10-16 and extraction by \texttt{arxiv\_latex}; these are supplied provenance notes, not a claim of independent access or verification. The registry attributes the collection to Albach, Barraud et al. and identifies Serge Tabachnikov as the contributor of this problem.
\item Reference label [KK]: the supplied statement says that this problem and its analogue for smooth Riemannian manifolds are discussed there. No bibliographic identification or link for [KK] is supplied.
\item Reference labels [dV] and [dV2]: the supplied statement says that these works study length growth of wave fronts. No bibliographic identifications or links for these references are supplied.
\end{itemize}
\end{document}
